
\subsubsection{5.1 Registry Assembly (RQ1)}

The H-E1-v2 pipeline assembled the registry meeting all gate criteria: n\_analyzable = 4,493 (gate: >= 200) and n\_features\_with\_variance = 3 (gate: >= 3).

\textbf{Table 1: Feature Variance Comparison (H-E1 vs. H-E1-v2)}

\begin{table}[htbp]
\centering
\resizebox{\columnwidth}{!}{%
\begin{tabular}{llll}
\toprule
Feature & H-E1 Variance & H-E1-v2 Variance & Status \\
\midrule
`dedup\_documented` & 0.002889 & 0.01512 & Non-zero (both) \\
`perplexity\_filter\_documented` & 0.000000 & 0.000000 & Zero (both) \\
`domain\_composition\_documented` & 0.002889 & 0.02006 & Non-zero (both) \\
`decontamination\_documented` & 0.000000 & 0.00200 & Non-zero in H-E1-v2 only \\
\bottomrule
\end{tabular}
}
\end{table}

The \texttt{perplexity\textbackslash \_filter\textbackslash \_documented} feature remained zero-variance in both runs despite targeted sampling. Manual inspection of LLaMA-2, Mistral, and Qwen cards confirmed that labs describe this practice using alternative terminology such as ''quality filtering,'' ''CCNet filtering,'' or ''fastText quality filter'' rather than ''perplexity filtering.''

\begin{figure}[htbp]
\centering
\includegraphics[width=\columnwidth]{figures/dropout_funnel.png}
\caption{Registry construction pipeline showing data collection stages and dropout funnel.}
\label{fig:dropout_funnel}
\end{figure}

\textbf{Figure 1.} Registry construction pipeline: 4,497 raw leaderboard rows to 4,493 analyzable models with documentation features.

\begin{figure}[htbp]
\centering
\includegraphics[width=\columnwidth]{figures/doc_score_distribution.png}
\caption{Distribution of doc\_score across 4,493 models.}
\label{fig:doc_score_distribution}
\end{figure}

\textbf{Figure 2.} Distribution of doc\_score across 4,493 registry models. 96.5\% of models have doc\_score = 0.

\begin{figure}[htbp]
\centering
\includegraphics[width=\columnwidth]{figures/family_breakdown.png}
\caption{Model family breakdown in the registry.}
\label{fig:family_breakdown}
\end{figure}

\textbf{Figure 3.} Model family composition of the registry.

\subsubsection{5.2 Scale Baseline (RQ2)}

\textbf{Table 2: OLS Baseline -- Scale-and-Architecture Model}

\begin{table}[htbp]
\centering
\resizebox{\columnwidth}{!}{%
\begin{tabular}{ll}
\toprule
Estimand & Value \\
\midrule
R-squared\_baseline & 0.4247 \\
log(params) coefficient & Positive, statistically significant \\
log(tokens) coefficient & Positive, statistically significant \\
\bottomrule
\end{tabular}
}
\end{table}

The scale-and-architecture baseline explains 42.47\% of V2 benchmark variance. This is substantially lower than the 60-75\% that might be expected based on V1 benchmark data, consistent with the observation that V2 tasks (GPQA, MUSR, MATH Level 5) resist scale saturation. The remaining 57.53\% of unexplained variance leaves room for additional explanatory factors.

\subsubsection{5.3 Documentation Effect (RQ3)}

\textbf{Table 3: OLS Proposed -- Scale + Documentation Model}

\begin{table}[htbp]
\centering
\resizebox{\columnwidth}{!}{%
\begin{tabular}{ll}
\toprule
Parameter & Value \\
\midrule
doc\_score coefficient (beta\_docs) & -3.4520 \\
p-value (beta\_docs) & 1.145 x $10^-8$ \\
R-squared\_proposed & 0.4289 \\
delta\_R-squared & +0.0042 \\
Wald F-test (nested model comparison) & p = 1.145 x $10^-8$ \\
\bottomrule
\end{tabular}
}
\end{table}

The F-test for nested model comparison (baseline vs. proposed) is equivalent to the t-test for beta\_docs under OLS with a single added regressor; the Wald F-statistic equals the squared t-statistic. The reported p-value confirms that adding doc\_score provides a statistically significant improvement over the baseline model at p < 0.001.

The documentation coefficient beta\_docs = -3.45 is statistically significant and negative. Models with more documented curation practices score 3.45 points lower per additional documented feature after controlling for scale and architecture. However, the practical effect size is small: delta\_R-squared = +0.0042 represents 0.42 percentage points of additional variance explained -- a statistically detectable but substantively modest increment.

\begin{figure}[htbp]
\centering
\includegraphics[width=\columnwidth]{figures/feature_coverage.png}
\caption{Documentation feature coverage across model families.}
\label{fig:feature_coverage}
\end{figure}

\textbf{Figure 4.} Documentation coverage by feature and model family.

\begin{figure}[htbp]
\centering
\includegraphics[width=\columnwidth]{figures/benchmark_heatmap.png}
\caption{Benchmark score heatmap across model families.}
\label{fig:benchmark_heatmap}
\end{figure}

\textbf{Figure 5.} Benchmark score heatmap by architecture family and documentation status.

\subsubsection{5.4 Identifying the Size-Vintage Confound}

The negative coefficient is most plausibly explained by a size-vintage confound. Well-documented models (LLaMA-2, Mistral-7B, Qwen-7B, Pythia, OLMo) are predominantly 7B-70B parameter base models from the 2023-2024 era. High-scoring V2 models tend to be newer, larger, or instruction-tuned derivatives that inherit undocumented base model curation.

Three competing explanations for beta\_docs < 0, ordered by qualitative assessment of plausibility (formal tests distinguishing these mechanisms are not available in this study and are deferred to propensity-matched analysis):

\begin{enumerate}
\item \textbf{Size confound (most plausible):} Documented base models are predominantly in the 7B-13B range; larger undocumented models score higher due to scale.
\item \textbf{Temporal confound:} Well-documented models predate the V2 benchmark launch.
\item \textbf{Fine-tuned model dominance:} 96.5\% of models are fine-tuned derivatives with doc\_score = 0; many score well on instruction-following tasks unrelated to pretraining curation.
\end{enumerate}

The delta\_R-squared = +0.0042 indicates that doc\_score adds detectable information beyond scale, but the direction of the effect is confounded. Causal interpretation of beta\_docs is not warranted from this analysis.
